\section{Implications and Limitations}

\paragraph{Scale diversity versus replication.}
Theorem~\ref{thm:matched} applies to arbitrary placement of $m$ points inside $[-T,T]$. Replication shrinks noise as $m^{-1/2}$ but cannot reveal curvature beyond the observed span. Pilot budgets should therefore include genuinely separated log-scales when feasible.

\paragraph{Coverage is conditional on visible assumptions.}
No bounded pilot interval can exclude every slower component. \ScaleCert\ requires an exponent range, positivity, observation bands, and any structural-mismatch radius to be declared. A broad class may yield a wide or unbounded interval; this is an abstention, not an implementation failure. The public experiment demonstrates that in-range fit alone cannot calibrate mismatch beyond the range.

\paragraph{Scope.}
The theory covers monotone positive spectral mixtures and independent Gaussian pilot errors. Nonmonotone phase transitions, adaptively selected pilot runs, correlated errors, and learned structural classes remain open. The OLS proposition exactly connects population and trained risks under Gaussian design, but analogous debiasing for ridge, early-stopped SGD, and neural networks requires additional analysis. Our public evaluation is deliberately descriptive because the released curves do not provide independent standard errors.

\section{Conclusion}

A high in-range $R^2$ does not certify a scaling-law extrapolation. Exact linear-regression constructions show that below-noise spectral mass can change the asymptotic exponent and compute-to-target requirement, while floor--exponent confounding creates a quantitative log-span barrier. \ScaleCert\ turns explicit assumptions into a target interval, with a rigorous continuous-exponent correction and an exact bridge to finite-sample OLS. The practical lesson is not that every extrapolation must be wide, but that confidence should contract only after pilot data have ruled out the alternatives that matter for the decision.
